\lecture{6}{12 March.}{}
\section{Inverse Function Theorem}
Recall that if \(f: \mathbb{R} \to \mathbb{R} \) is invertible, i.e. \(f^{-1} : \mathbb{R} \to \mathbb{R} \) exists s.t. \(f \circ f^{-1} (x) = f^{-1} \circ f (x) = x\) for all \(x \in \mathbb{R} \), and \(f^{\prime} (x_0) \neq 0\) for some \(x_0 \in \mathbb{R} \), then \(f^{-1}\) is differentiable at \(f(x_0)\) and  
\[
    \left( f^{-1} \right)^{\prime} (f(x_0)) = \frac{1}{f^{\prime} (x_0)},
\]      
which can be proved by chain rule.

In fact, one can say something even when \(f^{\prime} \) is not invertible, as long as we know that \(f\) is continuously differentiable. If \(f^{\prime} (x_0) \neq 0\), then \(f^{\prime} (x_0)\) must be either strictly positive or strictly negative (since we assume \(f^{\prime} \) to be continuous by \autoref{thm: differentiable implies continuous}), which implies \(f^{\prime} (x)\) is strictly positive for \(x\) near \(x_0\), or strictly negative for \(x\) near \(x_0\), and thus one-to-one in such small interval. In either case, \(f\) will become invertible if we restrict the domain and codomain of \(f\) to be sufficiently close to \(x_0\) and to \(f(x_0)\), respectively.            

 Now if \(f: \mathbb{R} ^n \to \mathbb{R} ^n\) and \(f \in C^1\), i.e. \(f\) is continuously differentiable. Suppose 
\[
    f^{\prime} (x_0) = Df(x_0) = \left[ \frac{\partial f_i(x_0)}{\partial x_j}  \right] 
\] 
is invertible. Then, \(\exists \) a neighborhood \(U\) of \(x_0\) and \(V\) of \(f(x_0)\) s.t. \(f: U \to V\) is bijective and its inverse \(f^{-1} : V \to U\) is continuous and differentiable at \(f(x_0)\) with 
\[
    \left( f^{-1} \right)^{\prime} (f(x_0)) = \left[ f^{\prime} (x_0) \right]^{-1}. 
\]       
This is the Inverse Function Theorem, and we will prove it later. 

First, we prepare a lemma. 
\begin{lemma} \label{lm: inverse of linear map is linear}
    Let \(T: \mathbb{R} ^n \to \mathbb{R} ^n\) be an invertible linear transformation . Then \(T^{-1} : \mathbb{R} ^n \to \mathbb{R} ^n\) is also linear.  
\end{lemma}
\begin{proof}
    \todo{DIY}
\end{proof}

\begin{theorem}[The Inverse Function Theorem] \label{thm: The inverse function theorem}
    Let \(E\) be an open subset of \(\mathbb{R} ^n\) and let \(f : E \to \mathbb{R} ^n\) be a function which is continuously differentiable on \(E\). Suppose \(x_0 \in E\) is the point s.t. the linear map \(f^{\prime}(x_0) : \mathbb{R} ^n \to \mathbb{R} ^n \) is invertible, then there exists an open set \(U \subseteq E\) s.t. \(x_0 \in U\) and an open set \(V \subseteq \mathbb{R} ^n\) s.t. \(f(x_0) \in V\) so that \(f: U \to V\) is a bijection. In particular, there is an inverse map \(f^{-1} : V \to U \). Furthermore, this inverse map is differentiable at \(f(x_0)\), and 
    \[
        \left( f^{-1} \right)^{\prime}  (f(x_0)) = \left( f^{\prime} (x_0) \right)^{-1}.
    \]              
    \begin{center}
\begin{tikzpicture}[x=1cm,y=1cm,>=Latex]
  \draw[gray!65,thin,-{Latex[length=2.7mm]}] (-7.8,0) -- (-2.8,0);
  \draw[gray!65,thin,-{Latex[length=2.7mm]}] (-5.3,-2.95) -- (-5.3,3.15);
  \draw[gray!65,thin,-{Latex[length=2.7mm]}] (2.8,0) -- (7.8,0);
  \draw[gray!65,thin,-{Latex[length=2.7mm]}] (5.4,-2.95) -- (5.4,3.15);

  \node[font=\bfseries\Large] at (-5.25,3.5) {$\mathbb{R}^n$};
  \node at (-4.1,3.5) {(domain)};
  \node[font=\bfseries\Large] at (5.35,3.5) {$\mathbb{R}^n$};
  \node at (6.65,3.5) {(codomain)};

  \fill[blue!15,draw=black,thick]
    plot[smooth cycle,tension=1.05]
    coordinates {(-7.1,-1.35) (-6.75,1.45) (-5.3,1.8) (-3.95,1.35) (-3.8,-1.05) (-4.8,-2.05) (-6.25,-2.0)};
  \fill[green!18,draw=black,thick]
    plot[smooth cycle,tension=1.05]
    coordinates {(3.8,-1.25) (3.95,1.45) (5.4,1.8) (6.7,1.45) (6.95,-0.95) (6.0,-1.95) (4.65,-2.0)};

  \node[font=\Large] at (-6.65,2.05) {$U$};
  \node[font=\Large] at (3.95,2.05) {$V$};

  \fill (-5.55,0.4) circle (0.08);
  \fill (-6.2,-1.15) circle (0.08);
  \fill (5.0,0.48) circle (0.08);
  \fill (6.0,-1.0) circle (0.08);

  \node at (-5.2,0.7) {$x_0$};
  \node at (-6.3,-1.55) {$x$};
  \node at (5.95,0.58) {$y_0=f(x_0)$};
  \node at (6,-1.55) {$f(x)$};

  \draw[thick,-{Latex[length=3mm]}] (-3.95,1.5) .. controls (-0.45,2.95) and (1.95,2.95) .. (3.9,1.55);
  \node at (0.0,3.1) {$f$};

  \draw[thick,-{Latex[length=3mm]}] (-5.48,0.45) .. controls (-2.4,1.4) and (1.6,1.4) .. (4.9,0.55);
  \draw[thick,-{Latex[length=3mm]}] (-6.1,-1.1) .. controls (-2.1,-0.1) and (2.0,-0.1) .. (5.9,-0.95);

  \draw[thick,-{Latex[length=3mm]}] (3.9,-1.95) .. controls (1.25,-3.0) and (-1.65,-3.0) .. (-3.95,-1.95);
  \node at (0.0,-3) {$f^{-1}$};

  \node at (-0.2,-1.55) {$(f^{-1})'(y_0)=(f'(x_0))^{-1}$};
  \node at (-5.3,-3.05) {$f'(x_0)$ invertible};
\end{tikzpicture}
\end{center}
\end{theorem}

\begin{proof}
    Suppose such \(f^{-1}\) exists, i.e. \(\exists U, V\) s.t. \(f: U \to V\) and \(f^{-1} : V \to U\) are bijective where \(x_0 \in U\) and \(f(x_0) \in V\). Since \(f^{\prime} (x_0)\) is invertible and \(f^{-1} (f(x_0))\) exists, we know 
    \[
        f^{-1} (f(x)) = x \text{ for } x \in U \text{ and } \left( f^{-1} \right)^{\prime} (f(x_0)) \circ (f^{\prime} (x_0)) = Id \text{ by chain rule.}    
    \] 

    Now we have to show such \(U, V\) exists. Let \(A = f^{\prime} (x_0)\), then 
    \[
        \lim_{x \to 0} \frac{\left\lVert f(x_0 + x) - f(x) - A(x) \right\rVert }{\lVert x \rVert } = 0.
    \]      
    Note that \(A^{-1}\) exists and 
    \begin{align*}
        A^{-1} \left( f(x_0 + x) - f(x_0) - Ax \right) &= A^{-1} (f(x_0 + x) - f(x)) - A^{-1} Ax \\
        &= A^{-1} \left( f(x_0 + x) - f(x_0) \right) - x.   
    \end{align*}   
    Note that 
    \begin{align*}
        \lim_{x \to 0} \frac{\left\lVert A^{-1} (f(x_0 + x) - f(x_0)) - x\right\rVert }{\lVert x \rVert } = 0 
    \end{align*}
    since 
    \begin{align*}
       0 &\le \lim_{x \to 0} \frac{\left\lVert A^{-1} (f(x_0 + x) - f(x_0)) - x\right\rVert }{\lVert x \rVert } = \lim_{x \to 0} \frac{  \left\lVert A^{-1} \left( f(x_0 + x) - f(x) - A(x) \right)  \right\rVert  }{\lVert x \rVert } \\
       &\le \lim_{x \to 0} \frac{\lVert A^{-1} \rVert \left\lVert f(x_0 + x) - f(x) - A(x) \right\rVert  }{\lVert x \rVert } = 0. 
    \end{align*}
    Let \(g(x) = A^{-1} (f(x_0 + x) - f(x_0)) - x\). Thus, 
    \[
        \lim_{x \to 0} \frac{\lVert g(x) \rVert }{\lVert x \rVert } = 0. 
    \]
    Note that \(g(0) = 0\). Thus, 
    \[
        \lim_{x \to 0} \frac{\lVert g(x) - g(0) - 0 \rVert }{\lVert x \rVert } = 0, 
    \] which shows \(g\) is differentiable at \(0\) with derivative \(g^{\prime} (0) = 0\). Now let's analyze \(g(x)\). Since \(f\) is continuously differentiable, \(g\) is also continuously differentiable. Since \(E\) is open and \(x_0 \in E\), there exists \(r_0 > 0\) s.t. \(B(x_0, r_0) \subseteq E\). Hence, \(g\) is defined on \(B(0, r_0)\). (For \(x \in B(0, r_0)\), \(\lVert x \rVert < r_0 \), so \(\lVert (x + x_0) - x_0 \rVert < r_0 \), i.e. \(x + x_0 \in B(x_0, r_0) \subseteq E\), so \(g(x)\) is well-defined.)      

    Thus, for all \(i, j\), \(\exists r > 0\) s.t. 
    \[
        \left\vert \frac{\partial g_j}{\partial x_i}  \right\vert \le \frac{1}{2\sqrt{n^5}} \text{ for } x \in B(0, r)  
    \]     
    since \(\frac{\partial g_j(0)}{\partial x_i} = 0\) (since \(g^{\prime} (0) = 0\) and thus \(g\)'s Jacobian matrix at \(0\) is \(0\) matrix) and \(\frac{\partial g_j}{\partial x_i} \) is continuous,  and hence
    \[
        \left\vert \frac{\partial g_j(x)}{\partial x_i} - \frac{\partial g_j(0)}{\partial x_i}   \right\vert\le \frac{1}{2\sqrt{n^5} } \text{ if } \lVert x \rVert < r \text{ for some small enough } r.   
    \] 
    By this, we can choose \(r\) small enough so that \(g\) is defined on \(B(0, r)\) and  
    \[
        \left\lVert \frac{\partial g(x)}{\partial x_j}  \right\rVert = \left( \sum_{i=1}^n \left( \frac{\partial g_i(x)}{\partial x_j}  \right)^2   \right)^{\frac{1}{2}}  \le \frac{1}{2n^2} \text{ for all } x \in B(0, r), \quad j = 1,2,\dots   
    \]
    Now if \(x, y \in B(0, r)\), then \(x + t(y - x) \in B(0, r)\) if \(0 \le t \le 1\). Also, 
    \begin{align*}
        g(y) - g(x) &= \int _0^1 \left( \frac{\mathrm{d}}{\mathrm{d}t} \left[ g(x + t (y - x)) \right]   \right) \, \mathrm{d} t = \int _0^1 \left[ Dg(x + t(y - x)) \circ (y - x) \right] \, \mathrm{d} t \\
        &= \int _0^1 \left( D_{(y-x)} g(x + t(y - x)) \right) \, \mathrm{d} t.     
    \end{align*}
    Now we estimate \(D_v g(x)\) for \(v = (v_1, v_2, \dots , v_n) \in \mathbb{R} ^n\). 
    \begin{align*}
        \lVert D_v g(x) \rVert &= \left\lVert \sum_{j=1}^n v_j \frac{\partial g(x)}{\partial x_j}   \right\rVert \le \sum_{j=1}^n \lVert v_j \rVert \left\lVert \frac{\partial g(x)}{\partial x_j}  \right\rVert \le \sum_{j=1}^n \lVert v \rVert \frac{1}{2n^2} = \frac{\lVert v \rVert }{2n} \le \frac{\lVert v \rVert }{2}.       
    \end{align*}    
    Thus, 
    \[
        \left\lVert D_{(y-x)} g(x + t(y - x)) \right\rVert \le \frac{\lVert x - y \rVert }{2}. 
    \]
    Hence, we have 
    \[
        \left\lVert g(y) - g(x) \right\rVert = \left\lVert \int _0^1 \left( D_{(y-x)} g(x + t (y - x)) \right) \, \mathrm{d} t   \right\rVert \le \frac{\lVert x - y \rVert }{2}.  
    \]    
    Recall \(A^{-1} (f(x_0 + x) - f(x)) = g(x) + x\) and \(g(0) = 0\), and we have shown
    \[
        \lVert g(x) - g(y) \rVert \le \frac{\lVert x - y \rVert }{2} \text{ for } x, y \in B(0, r).  
    \]  
    Let \(\Phi (x) = x + g(x)\). By \autoref{lm: first lemma for inverse function theorem}, \(\Phi :B(0, r) \to \mathbb{R} ^n\) is one to one and \(B\left( 0, \frac{r}{2} \right) \subseteq \Phi \left( B(0, r) \right)  \). Since \(\Phi \) is one to one on \(B(0, r)\). Let \(W = \Phi ^{-1} \left( B\left(0, \frac{r}{2} \right)  \right) \), then \(\Phi : W \to B\left( 0, \frac{r}{2} \right) \) is one to one and onto. Now we know \(\Phi \) is continuous (since \(g\) is composition of differentiable and thus continuous map, so \(g\) is continuous, and thus \(\Phi \) is continuous), so \(W\) is open (preimage of open set is open for continuous map). Also, \(0 \in W \) since \(\Phi (0) = 0 \in B\left( 0, \frac{r}{2} \right) \). Now we claim \(\Phi ^{-1} : B \left( 0, \frac{r}{2} \right) \to W \) is continuous. 
    Note that for all \(x, y \in W\) 
    \begin{align*}
        \lVert \Phi (x) - \Phi (y) \rVert &= \lVert x - y + g(x) - g(y) \rVert \ge \lVert x - y \rVert - \lVert g(x) - g(y) \rVert \ge \lVert x - y \rVert - \frac{1}{2} \lVert x - y \rVert = \frac{\lVert x - y \rVert }{2}.     
    \end{align*}

    Let \(x = \Phi ^{-1}(u)\) and \(y = \Phi ^{-1}(v)\) for \(u, v \in B\left( 0, \frac{r}{2} \right) \), then 
    \[
        \left\lVert \Phi \left( \Phi ^{-1}(u) \right) - \Phi \left( \Phi ^{-1}(v) \right)   \right\rVert \ge \frac{1}{2} \lVert \Phi ^{-1}(u) - \Phi ^{-1}(v) \rVert.  
    \]                
    Thus,
    \[
        \lVert u - v \rVert \ge \frac{1}{2} \left\lVert \Phi ^{-1}(u) - \Phi ^{-1}(v) \right\rVert. 
    \]
    This gives
    \[
        \left\lVert \Phi ^{-1}(u) - \Phi ^{-1}(v) \right\rVert \le 2 \lVert u - v \rVert,
    \]
    and implies \(\Phi ^{-1}\) is continuous. Thus, \(\Phi :W \to B\left( 0, \frac{r}{2} \right) \) is a homomorphism, i.e. \(\Phi \) is continuous and \(\Phi ^{-1}\) is also continuous. Since \(\Phi (0) = 0\) and \(\Phi \) is one-to-one, \(\Phi ^{-1} (0) = 0\). Also, note that 
    \[
        \lim_{z \to 0} \frac{\left\lVert \Phi ^{-1}(z) - \Phi ^{-1}(0) - z \right\rVert }{\lVert z \rVert } = \lim_{z \to 0} \frac{\lVert \Phi ^{-1}(z) - z \rVert }{\lVert z \rVert }.  
    \]     
    Note that for \(z \neq 0\), if we let \(\Phi ^{-1}(z) = x\), then \(z = \Phi (x) = x + g(x)\), and \(x \neq 0\), so 
    \[
        \frac{\left\lVert \Phi ^{-1}(z) - z \right\rVert }{\lVert z \rVert } = \frac{\left\lVert x - (x + g(x)) \right\rVert }{\lVert z \rVert } = \frac{\lVert g(x) \rVert }{\lVert z \rVert } = \frac{\lVert g(x) \rVert }{\lVert x \rVert } \cdot \frac{\lVert x \rVert }{\lVert z \rVert }.
    \]    
    Hence, 
    \begin{align*}
        \lim_{z \to 0} \frac{\left\lVert \Phi ^{-1}(z) - z \right\rVert }{\lVert z \rVert } &= \lim_{z \to 0} \frac{\lVert g(x) \rVert }{\lVert x \rVert } \cdot \frac{\lVert x \rVert }{\lVert z \rVert } = \lim_{z \to 0} \frac{\lVert g(x) \rVert }{\lVert x \rVert } \cdot \frac{\left\lVert \Phi ^{-1}(z) \right\rVert }{\lVert z \rVert }   
    \end{align*}
    Now if \(z \to 0\), then since \(\Phi ^{-1}\) is continuous, so \(x = \Phi ^{-1}(z) \to \Phi ^{-1}(0) = 0\). Also, when \(z \to 0\), we know 
    \[
        \left\lVert \Phi ^{-1}(z) - \Phi ^{-1}(0) \right\rVert \le 2 \lVert z - 0 \rVert \implies \frac{\left\lVert \Phi ^{-1}(z) \right\rVert }{\lVert z \rVert } \le 2,
    \]    
    and since
    \[
        \lim_{z \to 0} \frac{\lVert g(x) \rVert }{\lVert x \rVert } = \lim_{x \to 0} \frac{\lVert g(x) \rVert }{\lVert x \rVert } = 0,  
    \]
    so by Squeeze theorem, 
    \[
       \lim_{z \to 0} \frac{\left\lVert \Phi ^{-1}(z) - z \right\rVert }{\lVert z \rVert } = 0. 
    \]
    This shows \(\left( \Phi ^{-1} \right)^ {\prime} (0) = I\). 
    Recall that 
    \[
        \Phi (x) = x + g(x) = A^{-1} (f(x + x_0) - f(x_0))
    \]
    and \(\Phi : W \to B\left( 0, \frac{r}{2} \right) \) is a homomorpism. Now set 
    \[
        U = x_0 + W, \quad V = f(x_0) + A\left( B\left( 0, \frac{r}{2} \right)  \right). 
    \] 
    Since \(W\) is open and \(A\) is a linear isomorphism, both \(U\) and \(V\) are open subsets of \(\mathbb{R} ^n\). Moreover, \(x_0 \in U\) and \(f(x_0) \in V\). We claim that \(f\) maps \(U\) bijectively to \(V\). Let \(u = x + x_0 \in U\), where \(x \in W\), then since \(\Phi (W) = B\left( 0, \frac{r}{2} \right) \), we have \(\Phi (x) \in B\left( 0,\frac{r}{2} \right) \), and thus 
    \[
        f(u) = f(x + x_0) = f(x_0) + A \Phi(x) \in f(x_0) + A\left( B\left( 0, \frac{r}{2} \right)  \right) = V.  
    \] 
    This shows \(f(U) \subseteq V\).        
    Conversely, let \(y \in V\), then 
    \[
        A^{-1} (y - f(x_0)) \in B\left( 0, \frac{r}{2} \right). 
    \]       
    Because \(\Phi : W \to B\left( 0,\frac{r}{2} \right) \) is surjective, there exists \(x \in W\) s.t. 
    \[
        \Phi (x) = A^{-1} (y - f(x_0)).
    \]  
    For this \(x\) we have 
    \[
        f(x + x_0) = f(x_0) + A \Phi (x) = y.
    \] 
    This shows \(V \subseteq f(U)\) and thus \(f(U) = V\), i.e. \(f:U \to V\) is surjective. Since \(\Phi \) is one to one, so the \(x\) above is unique. It follows that \(f\) is one to one on \(U\). Thus, \(f:U \to V\) is bijective.       

    Now for \(y \in V\) we know \(y = f(x + x_0)\) for some \(x \in W\), so \(x + x_0 = f^{-1}(y)\).  Recall 
    \[
        \Phi (x) = A^{-1} (f(x + x_0) - f(x_0)) = A^{-1} (y - f(x_0)).
    \]
    Hence,
    \[
        x = \Phi ^{-1} \left( A^{-1} (y - f(x_0)) \right) \implies f^{-1}(y) - x_0 = \Phi ^{-1} \left( A^{-1} (y - f(x_0)) \right). 
    \]
    Thus, 
    \[
        f^{-1}(y) =  \Phi ^{-1} \left( A^{-1} (y - f(x_0)) \right) + x_0.
    \]
    This shows
    \begin{align*}
        \left( f^{-1} \right)^{\prime} (f(x_0)) &= \left( \Phi ^{-1} \right)^{\prime}  \left( A^{-1} (f(x_0) - f(x_0)) \right) (A^{-1})^{\prime} (0) = (\Phi ^{-1})^{\prime} (0) \circ A^{-1} = IA^{-1} = (f^{\prime} (x_0))^{-1}   
    \end{align*}
    since \(x_0 \in U\) so \(f(x_0) \in V\) and every linear transformation is differentiable at every point where \(\left( A^{-1} \right)(x) = A^{-1} \) for all \(x\).
    Thus, we're done.  
\end{proof}

\begin{remark}
Since $\det f'(x_0) \neq 0$ and $\det$ is continuous, there exists a neighborhood $B$ of $x_0$ such that $f'(x)$ is invertible for all $x \in B$.

Applying the inverse function theorem at each point $x \in B$, we obtain neighborhoods $U_x$ such that $f^{-1}$ is differentiable on $f(U_x)$ and
\[
(f^{-1})'(f(x)) = (f'(x))^{-1}.
\]

Restrict to a sufficiently small neighborhood $U' \subset B$ of $x_0$ and set $V' = f(U')$. Then $f^{-1}$ is differentiable on $V'$.

Moreover, for $y \in V'$,
\[
(f^{-1})'(y) = (f'(f^{-1}(y)))^{-1}.
\]
Since $f^{-1}$ is continuous, $f'$ is continuous, and matrix inversion is continuous on the set of invertible matrices, it follows that $(f^{-1})'$ is continuous on $V'$.

Hence, $f^{-1}$ is continuously differentiable on $V'$.
\end{remark}